%% =====================================================================
%% DRAFT replacement for \subsection{Do gateway blind spots persist across hops?}
%% (main.tex lines ~321-325, label sec:persist).
%%
%% NOT YET INCLUDED in main.tex. Integrate after the Phase 3 pipeline lands
%% (pipeline.py: hop2_rewrite -> hop2_report -> variant_correlation ->
%%  slip_gen_qwen -> ... -> heterogeneity) and every \PENDING{} below is
%% replaced from results/numbers.json + results/fig_phase3_*.png.
%%
%% Status of numbers in this draft:
%%   FINAL   — slip-pool debrief (slip_report_debrief), H3(a) repass (1600/1600),
%%             H3(b) rho_rewrite 138/158 = 0.873 [0.813, 0.917], control 67/300,
%%             simulation re-run with rho_rewrite (sim_final rc=0, sim_slip rc=0),
%%             H4 six-variant correlation + rotation (variant_correlation.json,
%%             fixed 2026-10-29: v5 scored on its stratified n=400 sample)
%%   FINAL   — qwen replication + beta-binomial ICC (heterogeneity.json,
%%             2026-10-29: ICC_secret 0.115/0.126, cross-model r=0.53 rho=0.34)
%%   FINAL   — Phase 2 ensemble sensitivity on the 616/1391 residue sentence
%%             (residue_ensemble.json: ensemble 275/1391 = 0.198 [0.178, 0.219],
%%              D2 75/1391 = 0.054, D3 388/1391 = 0.279, 0 undetermined)
%% =====================================================================

\subsection{Do gateway blind spots persist across hops?}
\label{sec:persist}

The relay estimate in the original version of this section rested on 20
relays (0.95) and was therefore a sensitivity analysis rather than a
measurement. We replaced it with a purpose-built boundary study: for each of
the 20 semantic development secrets we generated 160 notes (4 prompting
variants~$\times$~20 trials) with llama3.2:3b, giving 1\,600 notes that were
each put through the frozen gateway (lexical 0.60 OR cosine $\geq$ 0.78).
Of these, 209 passed the gateway while carrying the secret: an overall slip
rate of 0.131 (95\% Clopper--Pearson CI 0.114 to 0.148).

\paragraph{Slip rate depends on how the secret is phrased, not only on the model}
Across prompting variants the slip rate ranged from 0.207 (antonym,
0.169 to 0.251) through 0.165 (plain, 0.130 to 0.205) and 0.107
(paraphrase, 0.079 to 0.142) to 0.043 (compressed, 0.025 to 0.067):
compression, which drops the sentences a semantic filter would match,
is caught far more often than an antonymic rewrite of the same content.
Across secrets the spread was wider still, from 0.013 (three of the
20 secrets, 1/80 each) to 0.488 (\texttt{per-02}) and 0.450
(\texttt{hr-04}) --- a factor of nearly 40 between the easiest and the
hardest secret under one fixed gateway (Figure~\ref{fig:slips}).

\paragraph{A deterministic re-pass cannot repair a slip}
If a note slipped once, pushing the \emph{post-gateway} text through the
same gateway again changes nothing: 158 of 158 survived slips
(1.000, 95\% Wilson CI 0.976 to 1.000) still leaked after re-application,
and the same check on all 1\,600 notes found zero disagreements between
the original verdict and the re-pass verdict. This is expected --- the
gateway is a deterministic function of the text --- but it establishes
that blind spots are properties of the \emph{gateway--secret pair}, not
of a particular sampling seed: whatever slips past a fixed gateway
slips past it every time it is applied.

\paragraph{Redaction of a caught note does not remove the secret}
Of the 1\,391 notes the gateway did catch, 616 (0.443, 95\% Wilson CI
0.417 to 0.469) still carried detectable secret content under the
lexical detector D1 after redaction: the gateway redacts the sentence
it flags, but the same fact reappears in a header, a bullet list, or a
later paraphrase of the same sentence. Because D1 has content precision
0.71 against human labels (Section~\ref{sec:detector}), we re-scored
all 1\,391 redacted notes with the three-detector majority-vote
ensemble (D1 lexical, D2 NLI entailment, D3 cross-family LLM judge):
275 of 1\,391 notes still leaked under the ensemble (0.198, 95\% Wilson
CI 0.178 to 0.219; D2 alone 75/1\,391 = 0.054, D3 alone 388/1\,391 =
0.279, zero undetermined). Catching a note and sanitizing it are
therefore different events even under the stricter ensemble detector,
and only the first is what the gateway's catch rate measures.

\paragraph{Persistence under a next-agent rewrite}
To test whether a survived slip survives a genuine rewrite rather than
only a deterministic re-application, each of the 158 survived slips was
rewritten by an independent next agent (llama3.2:3b) and the result put
through the frozen gateway again.  138 of the 158 still carried the
secret: $\rho_{\text{rewrite}} = 0.873$ (95\% Wilson CI 0.813 to
0.917).  The matched control --- 300 notes caught at hop~1, rewritten
the same way --- originated the secret fresh in 67 cases, 0.223 (0.180
to 0.274).  A note that has already slipped is thus about four times
as likely to leak again after a full rewrite as a caught note is to
leak for the first time: blind spots persist across independent
rephrasing, at 0.87 rather than the 0.95 implied by the original
20-relay sample.

\paragraph{Do different gateways miss the same notes?}
Six gateways with deliberately different decision logic were scored on
every note: the frozen study gateway (lexical $0.60$ OR cosine
$\geq 0.78$), its lexical-only and cosine-only ablations, a
threshold-shifted variant (lexical $0.75$ OR cosine $\geq 0.85$), and a
secret-blind PII regex; the secret-blind regex+LLM-judge gateway was
additionally evaluated on a stratified sample of 400 notes (200 slips,
200 catches), so its rate is conditional on that sample.  Miss rates
(the gate passes while the note really leaks) were 0.125 for the frozen
gateway (95\% CI 0.110 to 0.142), 0.157 for cosine-only, 0.329 for the
threshold-shifted variant and 0.379 for lexical-only; the plain
secret-blind regex caught none of the 1\,600 notes (miss rate 0.994),
while regex+judge missed 0.145 of its 400-note sample.  Misses were
strongly correlated among variants that share decision components ---
$\phi$ = 0.876 frozen versus cosine-only, 0.539 frozen versus
threshold-shifted, 0.483 frozen versus lexical-only --- but negligible
between the study gateway and the secret-blind judge ($\phi$ = 0.061,
odds ratio 1.42; the other three deterministic judge pairs $|\phi|
\leq 0.063$, odds ratios 0.98 to 1.43), although every pair's marginal
miss rates differ (exact McNemar $p < 10^{-4}$ for all 15 pairs;
Figure~\ref{fig:h4}).  Correlation therefore tracks how much decision
logic two gateways share: genuinely different evidence finds genuinely
different blind spots, so H4's positive-correlation claim holds within a
decision family and is negligible across families.  The rotation claim
follows.  Re-applying the same gateway is deterministic (158/158
survived, above), and rotating within the family cannot help --- every
frozen miss is also a miss for lexical-only, cosine-only, the
threshold-shifted variant and the plain regex by construction, so
rotation stays at 1.000 --- whereas rotating to the secret-blind judge
cuts hop-2 persistence to 0.168 ($n$ = 400, an 83\% reduction).  H4's
rotation claim is thus supported only across genuinely different
decision logic.

\paragraph{How much of the variation is between secrets and between models?}
A beta-binomial fit over the 20 development secrets assigns a modest
but consistent share of the slip-rate variation to the secret itself:
$\mathrm{ICC}_{\text{secret}} = 0.115$ for llama3.2:3b and 0.126 for
the independent qwen2.5:3b replication pool (leave-one-secret-out mean
absolute prediction error 0.106 and 0.068), with the surviving-slip
arm giving 0.108 and 0.141.  The replication pool also measures the
model effect under the identical gateway and prompting design: 101 of
1\,600 notes slipped with qwen2.5:3b, 0.063 (95\% CI 0.052 to 0.076),
against 209 with llama3.2:3b, 0.131 (0.114 to 0.148) --- roughly half
the rate.  Which secrets are hard is only partly stable across
models: per-secret rates correlate moderately (Pearson $r = 0.53$,
Spearman $\rho = 0.34$ over 20 secrets), the two hardest secrets for
llama (\texttt{per-02} 0.488, \texttt{hr-04} 0.450) remain among the
hardest for qwen (0.213 each) while \texttt{cred-03} collapses from
0.375 to 0.038, and 7 of the 20 secrets never slip at all under qwen
versus none under llama.  Between-model dispersion is 0.322 on the
logit scale, but with $n_{\text{models}} = 2$ this is descriptive, not
a fitted random effect.  These measurements replace the ``0 to 0.10
per model'' origination range quoted from the three-agent runs.

\paragraph{Simulated outbreaks with a measured relay rate}
We re-ran the 500-agent simulations with
$\rho_{\text{rewrite}} = 0.873$ in place of the 0.95 relay estimate
(200 runs per point, as in the original study).  Under per-edge slip,
the outbreak floor with a gateway on every agent is 0.02 to 0.09
across the 18 model--network--placement cells (independent blocking
0.00; correlated blocking 0.04 to 0.09; pooled 2 to 9\% across models
and topologies).  As a sensitivity check using only the 40-episode
six-agent origination rates per model, the full-coverage floors are
0.00 for gemma2:2b, 0.09 to 0.12 for llama3.2:3b and 0.27 to 0.38 for
qwen2.5:3b, and the protected fraction needed to push outbreak below
5\% remains unreachable for the latter two models.  The floor is not
sensitive to the relay rate over this range: a controlled re-run at
the legacy $\rho = 0.95$ reproduces the previously published
0.09--0.14 and 0.24--0.32 per-model figures exactly, and a
1{,}000-run comparison of the two rates differs by at most 0.02 in
every cell, so the cell-level movement between runs is Monte Carlo
noise at the study's run count.  What changed is the evidential base:
the floor now rests on 158 measured relays instead of 20.

If blind spots persist like this, protected agents do not block
independently, and a cascade has two stages: a rare slip (origination
through the gateway) followed by near-certain relay. A back-of-envelope
check agrees: with $T = 0.017$ (95\% CI 0.006 to 0.031) and an origin
with four neighbours, the chance that at least one first-hop leak slips
through is about 7\% (2\% to 12\%). So the gateway cuts the probability
that a leak starts by about 40 times, but it does not make an outbreak
impossible once one slip occurs.

\begin{figure}[t]
\centering
\figbox{fig_phase3_slips.png}{0.85}
\caption{Phase 3 boundary pool (DEV, 1\,600 notes, llama3.2:3b). Left:
slip rate by prompting variant with exact 95\% Clopper--Pearson
intervals (dashed: pooled). Right: per-secret slip rate, worst first.}
\label{fig:slips}
\end{figure}

\begin{figure}[t]
\centering
\figbox{fig_phase3_persistence.png}{0.85}
\caption{Phase 3 persistence. Left: deterministic re-pass outcome for
survived, gate-passed and caught notes. Right: rewrite persistence
($\rho$) against the fresh-origination control, Wilson 95\% intervals.}
\label{fig:persist}
\end{figure}

\begin{figure}[t]
\centering
\figbox{fig_phase3_h4.png}{0.95}
\caption{Phase 3 H4: pairwise $\phi$ of miss indicators across six
gateway variants (blue: exact McNemar $p < 0.05$; all 15 pairs).
Study-gateway pairs use all 1\,600 notes; pairs involving the
secret-blind judge use its stratified 400-note sample.}
\label{fig:h4}
\end{figure}